\section{Decimal expansions of real numbers}
\label{sec:decimal-expansions}
A bar marks exactly the digits that repeat. Thus
$0.0\overline{2}=0.02222\ldots$, whereas
$0.\overline{02}=0.020202\ldots$.
A terminating decimal is understood to continue with zeros.
For a decimal $K.d_1d_2\cdots$, its value means the convergent series
\[
 K+\sum_{j=1}^{\infty}d_j10^{-j},\qquad 0\le d_j\le9.
\]
Convergence follows by comparison with the geometric series
$\sum_{j=1}^{\infty}9\cdot10^{-j}=1$.
Theorem~16.3 describes the only possible ambiguity: a positive
terminating decimal also has a representation ending in nines.

\exercise{16.1}{Two expansions for eleven fourths}
\goal Verify both decimal representations in part (a), and identify
which one long division produces in part (b).

\step{(a) Evaluate the expansion ending in zeros.}
The expression $2.75\overline{0}$ means
\[
 2+\frac7{10}+\frac5{100}
 =\frac{200+70+5}{100}
 =\frac{275}{100}=\frac{11}{4}.
\]
Every remaining term is zero, so there is no additional contribution.

\step{(a) Evaluate the expansion ending in nines.}
In $2.74\overline{9}$ the first nine is in the third decimal place.
The sum of the nines is
\[
 \sum_{j=3}^{\infty}9\cdot10^{-j}
 =9\cdot10^{-3}\sum_{k=0}^{\infty}(10^{-1})^k
 =\frac{9\cdot10^{-3}}{1-10^{-1}}
 =\frac1{100}.
\]
Consequently
\[
 2.74\overline{9}
 =2+\frac7{10}+\frac4{100}+\frac1{100}
 =2+\frac7{10}+\frac5{100}
 =\frac{11}{4}.
\]
The extra tail of nines supplies exactly the missing hundredth.

\step{(b) Carry out the long division.}
First $11=4\cdot2+3$, so the integer part is $2$ and the initial
remainder is $3$. Multiplying the remainder by ten gives
\[
 30=4\cdot7+2,\qquad 20=4\cdot5+0.
\]
The successive decimal digits are therefore $7$ and $5$, after
which the remainder is zero. All subsequent divisions have
$10\cdot0=4\cdot0+0$, so all subsequent digits are zero.
\answer{Both $2.74\overline9$ and $2.75\overline0$ equal $11/4$.
The long division algorithm produces $2.75\overline0$.}

\exercise{16.2}{Trading a terminating tail for a tail of nines}
\goal Verify the two identities in the opening paragraph of the
proof of Theorem~16.3: the identity for a positive integer, and the
identity for a decimal with a last nonzero digit.

\step{First calculate a general tail of nines.}
For any integer $r\ge0$, a geometric series gives
\[
 \sum_{j=r+1}^{\infty}9\cdot10^{-j}
 =\frac{9\cdot10^{-(r+1)}}{1-10^{-1}}
 =10^{-r}.
\]
In particular, when $r=0$ this tail has value $1$.

\step{A positive integer has the asserted alternative representation.}
Let $K$ be a positive integer. Then $K-1$ is a nonnegative integer,
so it is an allowed integer part for a nonnegative decimal. The
tail formula with $r=0$ gives
\[
 (K-1).\overline9
 =(K-1)+\sum_{j=1}^{\infty}9\cdot10^{-j}
 =(K-1)+1=K=K.\overline0.
\]
These digit strings are distinct even though their values agree.

\step{A last nonzero digit can be reduced by one.}
Suppose
\[
 x=K.d_1\cdots d_{r-1}d_r\overline0,\qquad r\ge1,
 \qquad d_r\in\{1,\ldots,9\}.
\]
The condition $d_r\ne0$ makes $d_r-1$ an allowed digit. The value
of the proposed other expansion is
\[
 \begin{aligned}
 &K+\sum_{j=1}^{r-1}d_j10^{-j}
       +(d_r-1)10^{-r}+\sum_{j=r+1}^{\infty}9\cdot10^{-j}\\
 &\qquad=K+\sum_{j=1}^{r-1}d_j10^{-j}
       +(d_r-1)10^{-r}+10^{-r}\\
 &\qquad=K+\sum_{j=1}^{r-1}d_j10^{-j}+d_r10^{-r}=x.
 \end{aligned}
\]
Thus decreasing the last nonzero digit by one and replacing the
following zeros by nines preserves the represented number.

\step{Check the zero case separately.}
The first identity requires $K>0$, and the second requires a last
nonzero digit. Neither condition holds for $0.\overline0$.
Indeed, if a nonnegative decimal represents zero, its nonnegative
integer part must be zero. Every digit must then be zero too:
a positive digit $d_j$ would contribute the positive amount
$d_j10^{-j}$ to a sum of nonnegative terms. Hence zero has just
the expansion $0.\overline0$.
\answer{The two claimed alternative expansions follow from the
geometric tail identity. Every positive terminating decimal has
the indicated second expansion; zero is the exceptional case.}

\exercise{16.3}{Strict comparison of convergent nonnegative series}
\goal Suppose $\sum a_n$ and $\sum b_n$ converge, all their terms
are nonnegative, $a_n\le b_n$ for every $n$, and $a_m<b_m$ at
at least one index $m$. Prove the inequality of their sums is strict.

\idea Keep the positive gap at that one index when passing from
finite sums to infinite sums. A strict inequality at each partial
sum alone would not guarantee a strict inequality of limits.

\step{Name the sums and the fixed positive gap.}
Write
\[
 A=\sum_{n=1}^{\infty}a_n,\qquad
 B=\sum_{n=1}^{\infty}b_n,\qquad
 \delta=b_m-a_m>0.
\]
Both $A$ and $B$ are finite by hypothesis. For every $N\ge m$,
\[
 \sum_{n=1}^{N}b_n-\sum_{n=1}^{N}a_n
 =\sum_{n=1}^{N}(b_n-a_n)
 \ge b_m-a_m=\delta.
\]
The inequality holds because all the other differences are
nonnegative; deleting them cannot increase the sum.

\step{Take the limit with the gap still present.}
The difference rule for limits gives
\[
 B-A=\lim_{N\to\infty}
 \left(\sum_{n=1}^{N}b_n-\sum_{n=1}^{N}a_n\right)\ge\delta>0.
\]
Here the weak inequality is preserved by passage to finite limits,
and the final strict inequality comes from the fixed number
$\delta$, which does not shrink with $N$.
\answer{$\displaystyle\sum_{n=1}^{\infty}a_n
<\sum_{n=1}^{\infty}b_n$.}

\exercise{16.4}{Converting six repeating decimals to fractions}
\goal Express each given repeating decimal as a fraction of integers.
\idea Group the repeated block as an integer and use a geometric
series. A block of $r$ digits with integer value $D$, repeated
immediately after the decimal point, represents
\[
 \frac{D}{10^r}+\frac{D}{10^{2r}}+\cdots
 =\frac{D10^{-r}}{1-10^{-r}}=\frac{D}{10^r-1}.
\]
A leading zero counts as a digit of the block: the block $02$ has
length two and integer value $2$.

\step{(a) The decimal $0.\overline2$.}
Here the block has length one and value $2$, so
\[
 0.\overline2=\frac2{10-1}=\boxed{\frac29}.
\]

\step{(b) The decimal $0.0\overline2$.}
Only the digit $2$ repeats; the first decimal digit is the single
nonrepeating zero. Moving the decimal point one place right gives
$10x=0.\overline2=2/9$, where $x=0.0\overline2$. Therefore
\[
 x=\frac1{10}\frac29=\frac2{90}=\boxed{\frac1{45}}.
\]
This differs from part (c), where the zero repeats as part of the block.

\step{(c) The decimal $0.\overline{02}$.}
The repeated block has length two, so
\[
 0.\overline{02}=\frac2{10^2-1}=\boxed{\frac2{99}}.
\]

\step{(d) The decimal $3.\overline{14}$.}
Separate the integer part before evaluating the length-two block:
\[
 3.\overline{14}=3+\frac{14}{99}
 =\frac{297+14}{99}=\boxed{\frac{311}{99}}.
\]

\step{(e) The decimal $0.\overline{10}$.}
The block is $10$, of length two and integer value $10$. Hence
\[
 0.\overline{10}=\boxed{\frac{10}{99}}.
\]

\step{(f) The decimal $0.\overline{1492}$.}
This block has length four and integer value $1492$, giving
\[
 0.\overline{1492}=\frac{1492}{10^4-1}
 =\boxed{\frac{1492}{9999}}.
\]
\answer{The six fractions, in order, are
$2/9$, $1/45$, $2/99$, $311/99$, $10/99$, and $1492/9999$.}

\exercise{16.5}{Decimal expansions of six rational numbers}
\goal Find the decimal expansions of the six prescribed fractions.
For terminating values, include both expansions supplied by
Theorem~16.3.

\step{(a) Expand $1/8$.}
At each stage divide ten times the current remainder by $8$:
\[
 10=8\cdot1+2,\qquad
 20=8\cdot2+4,\qquad
 40=8\cdot5+0.
\]
The quotients give digits $1,2,5$. A zero remainder produces zeros
forever. Reducing the last nonzero digit by one gives the other
expansion, so
\[
 \boxed{\frac18=0.125\overline0=0.124\overline9}.
\]

\step{(b) Expand $1/16$.}
The successive divisions are
\[
 \begin{aligned}
 10&=16\cdot0+10,&100&=16\cdot6+4,\\
 40&=16\cdot2+8,&80&=16\cdot5+0.
 \end{aligned}
\]
The initial quotient zero matters: it is the first decimal digit.
The digits are $0,6,2,5$, followed by zeros. Thus
\[
 \boxed{\frac1{16}=0.0625\overline0=0.0624\overline9}.
\]

\step{(c) Expand $2/3$.}
The first division is $20=3\cdot6+2$. The new remainder is the
original remainder $2$, so the same digit $6$ and remainder $2$
recur at every step. Therefore
\[
 \boxed{\frac23=0.\overline6}.
\]
As a check, the geometric-series formula gives
$0.\overline6=6/9=2/3$.

\step{(d) Expand $7/9$.}
Here $70=9\cdot7+7$. The remainder $7$ returns immediately, so
every digit is $7$:
\[
 \boxed{\frac79=0.\overline7}.
\]

\step{(e) Expand $6/11$.}
The divisions $60=11\cdot5+5$ and $50=11\cdot4+6$ return to the
initial remainder $6$ after the two digits $5,4$. Hence
\[
 \boxed{\frac6{11}=0.\overline{54}}.
\]
The block formula confirms $54/99=6/11$.

\step{(f) Expand $22/7$.}
First $22=7\cdot3+1$, giving integer part $3$ and remainder $1$.
The next six divisions are
\[
 \begin{aligned}
 10&=7\cdot1+3,&30&=7\cdot4+2,\\
 20&=7\cdot2+6,&60&=7\cdot8+4,\\
 40&=7\cdot5+5,&50&=7\cdot7+1.
 \end{aligned}
\]
The remainder has returned to $1$, so the six-digit block repeats:
\[
 \boxed{\frac{22}{7}=3.\overline{142857}}.
\]

\step{Account for all possible expansions.}
The last four displayed expansions have neither an eventual tail
of zeros nor an eventual tail of nines. Theorem~16.3 says they
are unique. The first two are positive terminating decimals, so
the two representations displayed in each of (a) and (b) are
their only representations.
\answer{The expansions are the six boxed answers above, including
both representations in the two terminating cases.}

\exercise{16.6}{The cyclic decimal blocks for the sevenths}
\goal Find the expansions of $1/7,2/7,\ldots,6/7$ and explain
the pattern in their repeating blocks.

\step{Record the long division rule once for each possible nonzero remainder.}
For a current remainder $r$, the next digit $d$ and next remainder
$r'$ satisfy $10r=7d+r'$ with $0\le r'<7$.
The six possible nonzero remainders give
\[
\begin{array}{c|rrrrrr}
 r&1&2&3&4&5&6\\ \hline
 d&1&2&4&5&7&8\\
 r'&3&6&2&5&1&4
\end{array}
\]
For instance, the entry $r=3$ comes from $30=7\cdot4+2$,
so the next digit is $4$ and the next remainder is $2$.

\step{Follow the cycle starting from remainder one.}
The remainder rule traces
\[
 1\longmapsto3\longmapsto2\longmapsto6
 \longmapsto4\longmapsto5\longmapsto1.
\]
The digits produced on these six steps are respectively
$1,4,2,8,5,7$. All six remainders in the cycle are distinct,
so returning to the initial remainder takes exactly six steps.

\step{Start the same cycle at each of the other remainders.}
The fraction $a/7$ begins with remainder $a$, for $a=1,\ldots,6$.
Starting the displayed cycle at that remainder gives
\[
 \begin{array}{c|c|c}
 \text{Fraction}&\text{Remainders before the successive digits}
                &\text{Expansion}\\ \hline
 1/7&1,3,2,6,4,5&0.\overline{142857}\\
 2/7&2,6,4,5,1,3&0.\overline{285714}\\
 3/7&3,2,6,4,5,1&0.\overline{428571}\\
 4/7&4,5,1,3,2,6&0.\overline{571428}\\
 5/7&5,1,3,2,6,4&0.\overline{714285}\\
 6/7&6,4,5,1,3,2&0.\overline{857142}
 \end{array}
\]
Each row returns to its initial remainder after the six digits,
so the same block repeats forever. These expansions are unique
because none has an eventual tail of zeros or nines.
\answer{The six blocks are cyclic rotations of $142857$.
Long division uses the same cycle of six nonzero remainders,
with a different starting position for each fraction.}

\exercise{16.7}{The decimal formed by concatenating integers}
\goal Decide whether the number
\[
 c=0.1234567891011121314151617181920212223\cdots
\]
is rational. Here its digits are obtained by writing $1,2,3,\ldots$
in their ordinary decimal notation and concatenating those strings.

\idea Prove this digit string is not eventually periodic. It has
arbitrarily long strings of zeros arbitrarily far out, but it also
has nonzero digits arbitrarily far out. Both facts are needed:
a decimal ending in zeros would have the first property alone.

\step{Find long zero strings after any prescribed position.}
For every positive integer $m$, the integer $10^m$ appears in the
concatenation. Its ordinary decimal notation is the digit $1$
followed by $m$ zeros. As $m$ increases, the starting positions
of these integers go arbitrarily far out: before writing $10^m$,
we have already written the $10^m-1$ preceding positive integers,
each of which contributes at least one digit.
Thus, after any specified digit position and for any specified
length $p$, there is a string of at least $p$ consecutive zeros.
Choose $m\ge p$ large enough that the occurrence of $10^m$
starts after that position.

\step{Find nonzero digits after any prescribed position.}
The initial digit $1$ in each of those occurrences of $10^m$ is
nonzero, and these occurrences start arbitrarily far out.
Consequently the decimal does not eventually consist only of zeros.
It does not eventually consist only of nines either, since zero
strings also occur arbitrarily far out.

\step{Rule out an eventual repeating block.}
Suppose the digits were periodic with period $p\ge1$ after some
position $N$. If the repeating block contained only zeros, all
digits beyond $N$ would be zero, contradicting the preceding step.
Otherwise the block contains a nonzero digit. Every run of $p$
consecutive positions beyond $N$ passes through every position of
the repeating block once, so it must contain a nonzero digit.
But the first step supplies a run of $p$ zeros entirely beyond
$N$. This is a contradiction in both cases.

\step{Apply the decimal characterization of rational numbers.}
The given decimal is not eventually periodic. By Theorem~16.5,
a rational number has a repeating decimal expansion; by
Theorem~16.3 its only possible second expansion would have tails
of zeros and nines. Those tails have already been excluded here.
Therefore the displayed decimal cannot represent a rational number.
\answer{The concatenation decimal is irrational.}

\exercise{16.8}{Constructing a number outside a proposed list}
\goal Given any sequence $(s_n)$ in $(0,1)$, with selected
expansions $s_n=0.d_1^{(n)}d_2^{(n)}\cdots$, define
\[
 e_n=\begin{cases}
 6,&d_n^{(n)}\ne6,\\
 7,&d_n^{(n)}=6.
 \end{cases}
\]
Show that $0.e_1e_2\cdots$ represents a number $y\in(0,1)$
different from every $s_n$.

\step{Check that the new digit string represents a number.}
Every $e_n$ is either $6$ or $7$, hence an allowed decimal digit.
Define the partial sums $y_N=\sum_{n=1}^{N}e_n10^{-n}$.
They increase and are bounded above by the convergent geometric
series $\sum_{n=1}^{\infty}7\cdot10^{-n}=7/9$.
The monotone sequence theorem gives a finite limit $y$, which
by the definition of a decimal is the number $0.e_1e_2\cdots$.
Termwise comparison of the finite partial sums and passage to
limits give
\[
 \frac23=\sum_{n=1}^{\infty}6\cdot10^{-n}
 \le y\le\sum_{n=1}^{\infty}7\cdot10^{-n}=\frac79.
\]
Because $0<2/3\le7/9<1$, these bounds imply $y\in(0,1)$.

\step{Verify the diagonal difference at every index.}
Fix $n$. If $d_n^{(n)}\ne6$, the construction gives $e_n=6$,
so $e_n\ne d_n^{(n)}$. If $d_n^{(n)}=6$, it gives $e_n=7$,
so again $e_n\ne d_n^{(n)}$.
Thus the new expansion differs from the selected expansion of
$s_n$ at least at its $n$th decimal digit.

\step{Explain why different strings cannot secretly represent the same number.}
Different decimal strings can represent the same real number in
the zero-tail/nine-tail case of Theorem~16.3. That issue must be
checked rather than ignored. Every digit of the expansion of $y$
is $6$ or $7$, so this expansion ends neither in zeros nor in
nines. Theorem~16.3 therefore makes it the unique expansion of
$y$. If $y=s_n$, the selected expansion of $s_n$ would be a
second expansion of $y$ differing at digit $n$, which is impossible.
Hence $y\ne s_n$ for every $n$.

\step{Interpret the conclusion for a supposed enumeration.}
If a sequence listed every number in $(0,1)$, the number $y$
constructed from that sequence would belong to $(0,1)$ and yet
would be absent from the list. This contradiction proves that
no such listing exists.
\answer{The constructed $y$ lies in $(0,1)$ and differs from every
listed term. Thus $(0,1)$ is uncountable.}

\exercise{16.9}{Existence of Euler's constant}
\goal For
\[
 H_n=\sum_{k=1}^{n}\frac1k,\qquad
 \gamma_n=H_n-\log n
 =H_n-\int_1^n\frac{dt}{t},
\]
prove (a) $(\gamma_n)$ decreases, (b) $0<\gamma_n\le1$, and
(c) its limit exists and is finite. The logarithm here is natural.

\step{(a) Subtract two neighboring terms without discarding the logarithm.}
Since $H_{n+1}=H_n+1/(n+1)$, we have
\[
 \begin{aligned}
 \gamma_n-\gamma_{n+1}
 &=H_n-\log n-H_n-\frac1{n+1}+\log(n+1)\\
 &=\int_n^{n+1}\frac{dt}{t}-\frac1{n+1}\\
 &=\int_n^{n+1}\left(\frac1t-\frac1{n+1}\right)dt.
 \end{aligned}
\]
For $n\le t<n+1$ the integrand is positive, because $t<n+1$
implies $1/t>1/(n+1)$. It is continuous and positive on the
interior of this interval, so its integral is strictly positive.
For an explicit lower estimate, on $[n,n+1/2]$ the integrand
is at least $1/(n+1/2)-1/(n+1)>0$, and that interval has positive
length $1/2$. Hence $\gamma_n>\gamma_{n+1}$ for every $n$.

\step{(b) Bound the integral between two rectangle sums.}
For $t\in[k,k+1]$ with $k\ge1$, reciprocation gives
\[
 \frac1{k+1}\le\frac1t\le\frac1k.
\]
Integrating over an interval of length one preserves these bounds:
\[
 \frac1{k+1}\le\int_k^{k+1}\frac{dt}{t}\le\frac1k.
\]
For $n\ge2$, sum from $k=1$ to $k=n-1$ to obtain
\[
 H_n-1\le\int_1^n\frac{dt}{t}\le H_n-\frac1n.
\]
Subtract the middle expression from $H_n$. Reversal of the
inequalities under subtraction gives
\[
 \frac1n\le\gamma_n\le1.
\]
The same conclusion holds for $n=1$ directly:
$\gamma_1=1-\log1=1$. In particular $\gamma_n>0$ at every index.

\step{(c) Apply monotone convergence with a real lower bound.}
Part (a) makes $(\gamma_n)$ decreasing, and part (b) bounds it
below by zero. The monotone sequence theorem therefore gives a
finite real limit
\[
 \gamma=\lim_{n\to\infty}\gamma_n\in[0,1].
\]
The fact that the two pieces $H_n$ and $\log n$ both grow without
bound does not obstruct convergence of their difference. The
monotonicity and bounds control that difference directly.
\answer{The sequence is strictly decreasing, satisfies
$0<\gamma_n\le1$, and has a finite limit called Euler's constant.}

\exercise{16.10}{Irrationality of the indicated powers and roots of pi}
\goal Starting with the irrationality of $\pi^2$ established in
Example~6(b), determine what follows for $\pi$, $\sqrt\pi$, and
$\sqrt[3]\pi$. Also address $\pi^4$ and distinguish the additional
fact needed for that last conclusion.

\step{Recall the elementary closure property of rational powers.}
If $q=a/b$ is rational, with $a,b\in\Z$ and $b\ne0$, then for
every positive integer $m$,
\[
 q^m=\frac{a^m}{b^m}\in\Q,
\]
since the numerator and denominator are integers and $b^m\ne0$.
Consequently, if a positive integer power of a number is
irrational, the number itself must be irrational.

\step{Apply this property to $\pi$ and to its square root.}
If $\pi$ were rational, its square would be rational, contrary
to the given result. Therefore $\pi$ is irrational.
If $\sqrt\pi$ were rational, its square would be $\pi$ and would
be rational, giving the same contradiction. Equivalently, its
fourth power is $\pi^2$, which is already known to be irrational.
Thus $\sqrt\pi$ is irrational too.

\step{Apply the same argument to the cube root.}
If $\sqrt[3]\pi$ were rational, its cube $\pi$ would be rational.
Since $\pi$ is irrational, this is impossible. One can also apply
the closure property directly to the sixth power:
$(\sqrt[3]\pi)^6=\pi^2$. Hence $\sqrt[3]\pi$ is irrational.

\step{Explain what the premise alone says about $\pi^4$.}
The implication just used runs from irrationality of a power to
irrationality of its base. It cannot be reversed: the square of
an irrational number can be rational. For a concrete comparison,
let $x=\sqrt[4]{2}$. Then $x^2=\sqrt2$ is irrational, but
$x^4=2$ is rational. Thus the statement ``$x^2$ is irrational''
does not by itself imply ``$x^4$ is irrational.'' In particular,
the given irrationality of $\pi^2$ alone does not settle the
question about $\pi^4$ by this rational-power argument.

\step{Determine the actual status of $\pi^4$ using the stronger fact stated in the text.}
Example~6(c) additionally states the deeper theorem that $\pi$
is not algebraic. This means it is not a zero of any nonzero
polynomial with integer coefficients. That stated theorem does
settle the question: if $\pi^4=a/b$ were rational, where
$a,b\in\Z$ and $b\ne0$, then
\[
 b\pi^4-a=0.
\]
The polynomial $P(X)=bX^4-a$ has integer coefficients and is
nonzero because its leading coefficient $b$ is nonzero. The
displayed equation would make $\pi$ algebraic, contradicting
Example~6(c). Therefore $\pi^4$ is irrational.

This final step cites the stronger theorem expressly stated in
the chapter; it is not a consequence of the irrationality of
$\pi^2$ alone.
\answer{$\pi$, $\sqrt\pi$, and $\sqrt[3]\pi$ are irrational by
elementary rational-power arguments from $\pi^2$. The number
$\pi^4$ is also irrational, using the stronger nonalgebraicity
result stated in Example~6(c).}
